\section{Four signs for the pair \texorpdfstring{\((d_n,d_{n^2})\)}{of second differences}}
\label{sec:signs}

Write $d_n=\lambda_{n+1}-2\lambda_n+\lambda_{n-1}$ for the centered
second difference of the classical Riemann Li coefficients.
The simultaneous signs of $d_n$ and $d_{n^2}$ describe the local convexity
or concavity of the same sequence at two widely separated indices.

\begin{corollary}\label{cor:signs}
For each $(a,b)\in\{-1,1\}^2$, the set
\[
 \{n\ge1:\ ad_n>11/1000,\quad bd_{n^2}>11/1000\}
\]
has positive lower Banach density. For the corresponding $h$ in
Table~\ref{tab:seeds}, these inequalities also hold at $n=h+p-1$
for a set of primes $p$ of positive lower density relative to the primes.
The same set of primes can be used for all four rows.
\end{corollary}

\begin{table}[ht]
\centering
\begin{tabular}{rcc@{\qquad}rr}
\hline
$h$ & $a$ & $b$ & lower bound for $ad_h$ & lower bound for $bd_{h^2}$\\
\hline
891327 & $+$ & $+$ & 0.01161 & 0.01387\\
842728 & $+$ & $-$ & 0.01180 & 0.01349\\
602346 & $-$ & $+$ & 0.01165 & 0.01133\\
673696 & $-$ & $-$ & 0.01119 & 0.01148\\
\hline
\end{tabular}
\caption{Certified strict inequalities for the Riemann zeta function.
Every decimal in the two right columns denotes an exact rational number.}
\label{tab:seeds}
\end{table}

\begin{proof}
First take the inequalities in the table, verified below. The polynomials
$u$ and $2hu+u^2$ vanish at $u=0$. The simultaneous uniform returns from
Theorem~\ref{thm:curv} therefore give a set of $u$ of positive lower
Banach density for which both $d_{h+u}$ and $d_{h^2+2hu+u^2}$ are arbitrarily
close to their respective values at $h$ and $h^2$. Take the tolerance
smaller than the least margin in the chosen row and put $n=h+u$.
For the prime assertion, the polynomials $p-1$ and
$(p-1)(2h+p-1)$ have the common unit root $1$ modulo every positive
integer. Take the finite family consisting of $p-1$ and all four
polynomials $(p-1)(2h+p-1)$. Their common unit root $1$ gives
simultaneous returns with $u=p-1$ for every row.
\end{proof}

\paragraph{Verification of Table~\ref{tab:seeds}.}
We use forty conjugate pairs of zeros. Their contribution can be certified
without resolving every zero below a height. Under RH a pair with positive
ordinate $t$ contributes $w(t)\cos(n\theta(t))$ to $d_n$, where
\[
 w(t)=\frac{2}{1/4+t^2},\qquad
 \theta(t)=2\arctan\!\frac{1}{2t}.
\]
The total mass of all pairs is $C=2+\gamma-\log(4\pi)$, as in
Theorem~\ref{thm:curv}. If $t_1,\ldots,t_{40}$ are the certified
ordinates, put $W=C-\sum_j w(t_j)$. For either $e\in\{-1,1\}$,
\begin{equation}
 ed_n\ \ge\ \sum_{j=1}^{40}e\,w(t_j)\cos\bigl(n\theta(t_j)\bigr)-W.
 \label{eq:signedtail}
\end{equation}
Here $0.01022325<W<0.01022326$.

The verifier starts with forty decimal centers $s_0=1/2+it_0$,
$13<t_0<130$. It assigns each the disk of radius $r=10^{-18}$ and
checks that the ordinate intervals are disjoint. At both $s=s_0$ and
$s=s_0+10^{-6}$ it evaluates $\zeta(s)$ by Euler--Maclaurin summation,
with cutoff $N=256$ and $m=20$ Bernoulli corrections. The explicit formula is
\[
 \sum_{j=1}^{N-1} j^{-s} + \frac{N^{1-s}}{s-1} + \frac{N^{-s}}{2}
 + \sum_{k=1}^{m}\frac{B_{2k}}{(2k)!}(s)_{2k-1}N^{1-s-2k} + R.
\]
Here $(s)_j$ is a rising factorial. The Fourier series for the periodic
Bernoulli polynomial gives
\[
 |R|\le \frac{2\zeta(2m)}{(2\pi)^{2m}}\,
 |(s)_{2m}|\frac{N^{1-\Re s-2m}}{\Re s+2m-1}
 \le 4N\left(\frac{171}{6N}\right)^{40}<10^{-34}.
\]
The code generates the Bernoulli numbers over the rationals and performs
all transcendental evaluations with outward-rounded 192-bit MPFI intervals.

For every center, the calculation gives $|\zeta(s_0)|<2\cdot10^{-25}$
and a coordinate of the difference quotient with absolute value greater
than $0.1$. To turn this into a zero certificate, first-order Euler
summation on $|s-s_0|\le1/4$ gives
$|\zeta(s)|\le |s-1|^{-1}+1/2+|s|/(2\Re s)<300$.
By Cauchy's estimate the Taylor terms of degree at least two sum in
absolute value to at most
$300(4|v|)^2/(1-4|v|)<5000|v|^2$ for $0<|v|\le10^{-6}$.
Thus $|\zeta'(s_0)|>0.1-5000\cdot10^{-6}>0.09$.
On $|v|=r$ the linear term exceeds in modulus the sum of the constant
term and the Taylor remainder. Rouch\'e's theorem gives one zero
counted with multiplicity inside each disk. Inserting each entire
ordinate interval into \eqref{eq:signedtail} yields
Table~\ref{tab:seeds}. All seeds, enclosures, and comparisons can be
reproduced by running \texttt{certify\_zeta.py}.
